\documentclass[11pt]{article}
\newcommand{\videoid}{-dMbSPakLgw}
% Preámbulo compartido de las notas del curso de topología algebraica.
% Cada nota empieza con \documentclass[11pt]{article}, define \videoid{...}
% (el identificador del video de YouTube) y luego hace % Preámbulo compartido de las notas del curso de topología algebraica.
% Cada nota empieza con \documentclass[11pt]{article}, define \videoid{...}
% (el identificador del video de YouTube) y luego hace % Preámbulo compartido de las notas del curso de topología algebraica.
% Cada nota empieza con \documentclass[11pt]{article}, define \videoid{...}
% (el identificador del video de YouTube) y luego hace \input{preambulo}.
% Compilar con LuaLaTeX:  lualatex nombre-de-la-nota.tex
\usepackage[spanish,provide=*]{babel}
\usepackage[T1]{fontenc}
\usepackage{lmodern}
\usepackage[a4paper,margin=2.6cm]{geometry}
\usepackage{amsmath,amssymb,amsthm}
\usepackage{xcolor}
\usepackage{tikz}
\usetikzlibrary{arrows.meta,babel}
\usepackage[colorlinks=true,urlcolor=blue!50!black,linkcolor=blue!50!black]{hyperref}

\hyphenpenalty=10000 \exhyphenpenalty=10000 \emergencystretch=3em

% Colores de las figuras (los mismos que usa la página)
\definecolor{extremo}{RGB}{47,107,79}
\definecolor{azul}{RGB}{44,74,154}

% Definiciones numeradas por unidad
\theoremstyle{definition}
\newtheorem{defn}{Definición}
% Numeración por unidad: \unidad{u}{n} (después de \input{preambulo}) hace que las
% definiciones de la nota se llamen u.n, u.(n+1), ...; n es el primer número
% libre de esa unidad (ver la tabla de CONVENCIONES.md).
\newcommand{\unidad}[2]{\renewcommand{\thedefn}{#1.\arabic{defn}}\setcounter{defn}{\numexpr#2-1\relax}}

% Observaciones y erratas del video. Usar SOLO si hay algo que señalar:
%   \begin{erratas}
%     \errata{203}{3:23}{Se dice ..., pero debe decir ...}
%   \end{erratas}
% El primer argumento son los segundos (enlaza al minuto exacto del video),
% el segundo es la marca visible, el tercero es el texto.
\newenvironment{erratas}{\section*{Observaciones y erratas del video}\begin{itemize}}{\end{itemize}}
\newcommand{\errata}[3]{%
  \edef\errurl{https://www.youtube.com/watch?v=\videoid\string&t=#1}%
  \item[\textbf{\expandafter\href\expandafter{\errurl}{#2}}] #3}
.
% Compilar con LuaLaTeX:  lualatex nombre-de-la-nota.tex
\usepackage[spanish,provide=*]{babel}
\usepackage[T1]{fontenc}
\usepackage{lmodern}
\usepackage[a4paper,margin=2.6cm]{geometry}
\usepackage{amsmath,amssymb,amsthm}
\usepackage{xcolor}
\usepackage{tikz}
\usetikzlibrary{arrows.meta,babel}
\usepackage[colorlinks=true,urlcolor=blue!50!black,linkcolor=blue!50!black]{hyperref}

\hyphenpenalty=10000 \exhyphenpenalty=10000 \emergencystretch=3em

% Colores de las figuras (los mismos que usa la página)
\definecolor{extremo}{RGB}{47,107,79}
\definecolor{azul}{RGB}{44,74,154}

% Definiciones numeradas por unidad
\theoremstyle{definition}
\newtheorem{defn}{Definición}
% Numeración por unidad: \unidad{u}{n} (después de % Preámbulo compartido de las notas del curso de topología algebraica.
% Cada nota empieza con \documentclass[11pt]{article}, define \videoid{...}
% (el identificador del video de YouTube) y luego hace \input{preambulo}.
% Compilar con LuaLaTeX:  lualatex nombre-de-la-nota.tex
\usepackage[spanish,provide=*]{babel}
\usepackage[T1]{fontenc}
\usepackage{lmodern}
\usepackage[a4paper,margin=2.6cm]{geometry}
\usepackage{amsmath,amssymb,amsthm}
\usepackage{xcolor}
\usepackage{tikz}
\usetikzlibrary{arrows.meta,babel}
\usepackage[colorlinks=true,urlcolor=blue!50!black,linkcolor=blue!50!black]{hyperref}

\hyphenpenalty=10000 \exhyphenpenalty=10000 \emergencystretch=3em

% Colores de las figuras (los mismos que usa la página)
\definecolor{extremo}{RGB}{47,107,79}
\definecolor{azul}{RGB}{44,74,154}

% Definiciones numeradas por unidad
\theoremstyle{definition}
\newtheorem{defn}{Definición}
% Numeración por unidad: \unidad{u}{n} (después de \input{preambulo}) hace que las
% definiciones de la nota se llamen u.n, u.(n+1), ...; n es el primer número
% libre de esa unidad (ver la tabla de CONVENCIONES.md).
\newcommand{\unidad}[2]{\renewcommand{\thedefn}{#1.\arabic{defn}}\setcounter{defn}{\numexpr#2-1\relax}}

% Observaciones y erratas del video. Usar SOLO si hay algo que señalar:
%   \begin{erratas}
%     \errata{203}{3:23}{Se dice ..., pero debe decir ...}
%   \end{erratas}
% El primer argumento son los segundos (enlaza al minuto exacto del video),
% el segundo es la marca visible, el tercero es el texto.
\newenvironment{erratas}{\section*{Observaciones y erratas del video}\begin{itemize}}{\end{itemize}}
\newcommand{\errata}[3]{%
  \edef\errurl{https://www.youtube.com/watch?v=\videoid\string&t=#1}%
  \item[\textbf{\expandafter\href\expandafter{\errurl}{#2}}] #3}
) hace que las
% definiciones de la nota se llamen u.n, u.(n+1), ...; n es el primer número
% libre de esa unidad (ver la tabla de CONVENCIONES.md).
\newcommand{\unidad}[2]{\renewcommand{\thedefn}{#1.\arabic{defn}}\setcounter{defn}{\numexpr#2-1\relax}}

% Observaciones y erratas del video. Usar SOLO si hay algo que señalar:
%   \begin{erratas}
%     \errata{203}{3:23}{Se dice ..., pero debe decir ...}
%   \end{erratas}
% El primer argumento son los segundos (enlaza al minuto exacto del video),
% el segundo es la marca visible, el tercero es el texto.
\newenvironment{erratas}{\section*{Observaciones y erratas del video}\begin{itemize}}{\end{itemize}}
\newcommand{\errata}[3]{%
  \edef\errurl{https://www.youtube.com/watch?v=\videoid\string&t=#1}%
  \item[\textbf{\expandafter\href\expandafter{\errurl}{#2}}] #3}
.
% Compilar con LuaLaTeX:  lualatex nombre-de-la-nota.tex
\usepackage[spanish,provide=*]{babel}
\usepackage[T1]{fontenc}
\usepackage{lmodern}
\usepackage[a4paper,margin=2.6cm]{geometry}
\usepackage{amsmath,amssymb,amsthm}
\usepackage{xcolor}
\usepackage{tikz}
\usetikzlibrary{arrows.meta,babel}
\usepackage[colorlinks=true,urlcolor=blue!50!black,linkcolor=blue!50!black]{hyperref}

\hyphenpenalty=10000 \exhyphenpenalty=10000 \emergencystretch=3em

% Colores de las figuras (los mismos que usa la página)
\definecolor{extremo}{RGB}{47,107,79}
\definecolor{azul}{RGB}{44,74,154}

% Definiciones numeradas por unidad
\theoremstyle{definition}
\newtheorem{defn}{Definición}
% Numeración por unidad: \unidad{u}{n} (después de % Preámbulo compartido de las notas del curso de topología algebraica.
% Cada nota empieza con \documentclass[11pt]{article}, define \videoid{...}
% (el identificador del video de YouTube) y luego hace % Preámbulo compartido de las notas del curso de topología algebraica.
% Cada nota empieza con \documentclass[11pt]{article}, define \videoid{...}
% (el identificador del video de YouTube) y luego hace \input{preambulo}.
% Compilar con LuaLaTeX:  lualatex nombre-de-la-nota.tex
\usepackage[spanish,provide=*]{babel}
\usepackage[T1]{fontenc}
\usepackage{lmodern}
\usepackage[a4paper,margin=2.6cm]{geometry}
\usepackage{amsmath,amssymb,amsthm}
\usepackage{xcolor}
\usepackage{tikz}
\usetikzlibrary{arrows.meta,babel}
\usepackage[colorlinks=true,urlcolor=blue!50!black,linkcolor=blue!50!black]{hyperref}

\hyphenpenalty=10000 \exhyphenpenalty=10000 \emergencystretch=3em

% Colores de las figuras (los mismos que usa la página)
\definecolor{extremo}{RGB}{47,107,79}
\definecolor{azul}{RGB}{44,74,154}

% Definiciones numeradas por unidad
\theoremstyle{definition}
\newtheorem{defn}{Definición}
% Numeración por unidad: \unidad{u}{n} (después de \input{preambulo}) hace que las
% definiciones de la nota se llamen u.n, u.(n+1), ...; n es el primer número
% libre de esa unidad (ver la tabla de CONVENCIONES.md).
\newcommand{\unidad}[2]{\renewcommand{\thedefn}{#1.\arabic{defn}}\setcounter{defn}{\numexpr#2-1\relax}}

% Observaciones y erratas del video. Usar SOLO si hay algo que señalar:
%   \begin{erratas}
%     \errata{203}{3:23}{Se dice ..., pero debe decir ...}
%   \end{erratas}
% El primer argumento son los segundos (enlaza al minuto exacto del video),
% el segundo es la marca visible, el tercero es el texto.
\newenvironment{erratas}{\section*{Observaciones y erratas del video}\begin{itemize}}{\end{itemize}}
\newcommand{\errata}[3]{%
  \edef\errurl{https://www.youtube.com/watch?v=\videoid\string&t=#1}%
  \item[\textbf{\expandafter\href\expandafter{\errurl}{#2}}] #3}
.
% Compilar con LuaLaTeX:  lualatex nombre-de-la-nota.tex
\usepackage[spanish,provide=*]{babel}
\usepackage[T1]{fontenc}
\usepackage{lmodern}
\usepackage[a4paper,margin=2.6cm]{geometry}
\usepackage{amsmath,amssymb,amsthm}
\usepackage{xcolor}
\usepackage{tikz}
\usetikzlibrary{arrows.meta,babel}
\usepackage[colorlinks=true,urlcolor=blue!50!black,linkcolor=blue!50!black]{hyperref}

\hyphenpenalty=10000 \exhyphenpenalty=10000 \emergencystretch=3em

% Colores de las figuras (los mismos que usa la página)
\definecolor{extremo}{RGB}{47,107,79}
\definecolor{azul}{RGB}{44,74,154}

% Definiciones numeradas por unidad
\theoremstyle{definition}
\newtheorem{defn}{Definición}
% Numeración por unidad: \unidad{u}{n} (después de % Preámbulo compartido de las notas del curso de topología algebraica.
% Cada nota empieza con \documentclass[11pt]{article}, define \videoid{...}
% (el identificador del video de YouTube) y luego hace \input{preambulo}.
% Compilar con LuaLaTeX:  lualatex nombre-de-la-nota.tex
\usepackage[spanish,provide=*]{babel}
\usepackage[T1]{fontenc}
\usepackage{lmodern}
\usepackage[a4paper,margin=2.6cm]{geometry}
\usepackage{amsmath,amssymb,amsthm}
\usepackage{xcolor}
\usepackage{tikz}
\usetikzlibrary{arrows.meta,babel}
\usepackage[colorlinks=true,urlcolor=blue!50!black,linkcolor=blue!50!black]{hyperref}

\hyphenpenalty=10000 \exhyphenpenalty=10000 \emergencystretch=3em

% Colores de las figuras (los mismos que usa la página)
\definecolor{extremo}{RGB}{47,107,79}
\definecolor{azul}{RGB}{44,74,154}

% Definiciones numeradas por unidad
\theoremstyle{definition}
\newtheorem{defn}{Definición}
% Numeración por unidad: \unidad{u}{n} (después de \input{preambulo}) hace que las
% definiciones de la nota se llamen u.n, u.(n+1), ...; n es el primer número
% libre de esa unidad (ver la tabla de CONVENCIONES.md).
\newcommand{\unidad}[2]{\renewcommand{\thedefn}{#1.\arabic{defn}}\setcounter{defn}{\numexpr#2-1\relax}}

% Observaciones y erratas del video. Usar SOLO si hay algo que señalar:
%   \begin{erratas}
%     \errata{203}{3:23}{Se dice ..., pero debe decir ...}
%   \end{erratas}
% El primer argumento son los segundos (enlaza al minuto exacto del video),
% el segundo es la marca visible, el tercero es el texto.
\newenvironment{erratas}{\section*{Observaciones y erratas del video}\begin{itemize}}{\end{itemize}}
\newcommand{\errata}[3]{%
  \edef\errurl{https://www.youtube.com/watch?v=\videoid\string&t=#1}%
  \item[\textbf{\expandafter\href\expandafter{\errurl}{#2}}] #3}
) hace que las
% definiciones de la nota se llamen u.n, u.(n+1), ...; n es el primer número
% libre de esa unidad (ver la tabla de CONVENCIONES.md).
\newcommand{\unidad}[2]{\renewcommand{\thedefn}{#1.\arabic{defn}}\setcounter{defn}{\numexpr#2-1\relax}}

% Observaciones y erratas del video. Usar SOLO si hay algo que señalar:
%   \begin{erratas}
%     \errata{203}{3:23}{Se dice ..., pero debe decir ...}
%   \end{erratas}
% El primer argumento son los segundos (enlaza al minuto exacto del video),
% el segundo es la marca visible, el tercero es el texto.
\newenvironment{erratas}{\section*{Observaciones y erratas del video}\begin{itemize}}{\end{itemize}}
\newcommand{\errata}[3]{%
  \edef\errurl{https://www.youtube.com/watch?v=\videoid\string&t=#1}%
  \item[\textbf{\expandafter\href\expandafter{\errurl}{#2}}] #3}
) hace que las
% definiciones de la nota se llamen u.n, u.(n+1), ...; n es el primer número
% libre de esa unidad (ver la tabla de CONVENCIONES.md).
\newcommand{\unidad}[2]{\renewcommand{\thedefn}{#1.\arabic{defn}}\setcounter{defn}{\numexpr#2-1\relax}}

% Observaciones y erratas del video. Usar SOLO si hay algo que señalar:
%   \begin{erratas}
%     \errata{203}{3:23}{Se dice ..., pero debe decir ...}
%   \end{erratas}
% El primer argumento son los segundos (enlaza al minuto exacto del video),
% el segundo es la marca visible, el tercero es el texto.
\newenvironment{erratas}{\section*{Observaciones y erratas del video}\begin{itemize}}{\end{itemize}}
\newcommand{\errata}[3]{%
  \edef\errurl{https://www.youtube.com/watch?v=\videoid\string&t=#1}%
  \item[\textbf{\expandafter\href\expandafter{\errurl}{#2}}] #3}

\unidad{1}{18} % unidad 1; primer enunciado: 1.18

\title{El grupo fundamental del círculo (parte 2)}
\author{Luis Jorge Sánchez Saldaña\\[2pt]\small Facultad de Ciencias, UNAM}
\date{}

\begin{document}
\maketitle

\noindent\small Notas del video
\href{https://www.youtube.com/watch?v=-dMbSPakLgw}{\guillemotleft{}El grupo fundamental del círculo -- parte 2\guillemotright{}}
(curso de Introducción a la Topología Algebraica). Se termina el cálculo: $\pi_1(S^1)\cong\mathbb{Z}$.
\normalsize

\section*{Un corolario de los lemas de levantamiento}

Se recuerda lo visto en el video anterior. La función $\varphi\colon\mathbb{R}\to S^1$, $\varphi(x)=e^{2\pi i x}$, enrolla $\mathbb{R}$ sobre $S^1$, y todo camino de $S^1$ que comienza en $1$ se puede levantar a $\mathbb{R}$ comenzando en $0$ (Lema~1.16); además, cualquier homotopía relativa a los extremos entre dos caminos se puede levantar (Lema~1.17). Con esto se obtiene un corolario inmediato.

\begin{cor}
Sean $f,g\colon I\to S^1$ caminos con $f(0)=g(0)=1$ que son homotópicos relativamente a los extremos, $f\simeq g\ \mathrm{rel}\ \{0,1\}$. Entonces sus levantamientos $\tilde f,\tilde g\colon I\to\mathbb{R}$ (con punto inicial $0$) son homotópicos relativamente a los extremos. En particular, $\tilde f(1)=\tilde g(1)$.
\end{cor}

\noindent\textbf{Comentario.} En el pizarrón el corolario se enuncia para $f,g\colon I\to S^1$ homotópicos rel $\{0,1\}$; la condición $f(0)=g(0)=1$ es la de los lemas de levantamiento, y la segunda afirmación es la que el video dice \guillemotleft{}en particular\guillemotright{}.

\medskip
La idea es la siguiente. Se tienen dos caminos en $S^1$ con los mismos extremos y una deformación de uno en el otro que deja los extremos fijos. Al levantarlos (con punto inicial $0$), no se sabe cómo quedan, pero independientemente de eso los levantamientos son de nuevo homotópicos relativamente a los extremos. Como esa deformación deja los extremos fijos, ambos levantamientos tienen el mismo punto inicial, el $0$, y el mismo punto final, que no se sabe cuál es pero sigue siendo el mismo para ambos.

\section*{El teorema}

\begin{teo}[El grupo fundamental del círculo]
$\pi_1(S^1)\cong\mathbb{Z}$ como grupos abelianos. Es decir, el grupo fundamental del círculo es un grupo cíclico infinito.
\end{teo}

\begin{proof}
Se define una función del grupo fundamental del círculo en los enteros: a la clase de $f$ se le asigna el punto final de su levantamiento,
\[
  \Phi\colon\pi_1(S^1,1)\longrightarrow\mathbb{Z},
  \qquad
  [f]\longmapsto\tilde f(1).
\]
\textbf{Comentario.} En el pizarrón se usa la misma letra $\varphi$ para esta función y para la función exponencial $\mathbb{R}\to S^1$; aquí se escribe $\Phi$ para el homomorfismo, para no confundirlas. Además, $\tilde f(1)$ es en efecto un entero: $\varphi(\tilde f(1))=f(1)=1$ y $e^{2\pi i x}=1$ si y solo si $x\in\mathbb{Z}$.

Hay que demostrar que $\Phi$ está bien definida (los elementos del dominio son clases de equivalencia, así que hay que ver que no depende del representante), que es un homomorfismo y que es biyectiva.

\medskip\noindent\textit{$\Phi$ está bien definida.}
Si $[f]=[g]$, por definición de la relación de equivalencia $f\simeq g\ \mathrm{rel}\ \{0,1\}$. Por el Corolario~1.18, los puntos finales de los levantamientos son iguales:
\[
  \Phi[f]=\tilde f(1)=\tilde g(1)=\Phi[g].
\]

\medskip\noindent\textit{$\Phi$ es un homomorfismo.}
Sean $[f],[g]\in\pi_1(S^1,1)$ y sean $m=\tilde f(1)$ y $n=\tilde g(1)$ los puntos finales de sus levantamientos. Se define el camino $f'\colon I\to\mathbb{R}$ como $\tilde f$ desplazado $n$ unidades (a la derecha o a la izquierda según el signo de $n$):
\[
  f'(s)=\tilde f(s)+n.
\]
Entonces $\tilde g$ va de $0$ a $n$ y $f'$ va de $n$ a $n+m$ (figura), de modo que $\tilde g\cdot f'$ tiene sentido.

\begin{center}
\begin{tikzpicture}[x=0.021cm,y=-0.021cm,line cap=round]
  \draw[gray] (20,70) -- (460,70);
  \draw[very thick] (60,70) -- (220,70);
  \draw[very thick,azul] (220,70) -- (380,70);
  \fill[extremo] (60,70) circle (4);
  \fill (220,70) circle (4);
  \fill[azul] (380,70) circle (4);
  \node at (140,45) {$\tilde g$};
  \node[azul] at (300,45) {$f'$};
  \node at (60,96) {\small $0$};
  \node at (220,96) {\small $n$};
  \node at (380,96) {\small $n+m$};
  \node at (452,48) {$\mathbb{R}$};
\end{tikzpicture}
\end{center}

\noindent Figura: el camino $f'$ está diseñado para poder concatenarse con $\tilde g$.

\medskip
Se verifica fácilmente que $\tilde g\cdot f'$ es el levantamiento de $g\cdot f$, el producto de los dos lazos: comienza en $0$ y su composición con $\varphi$ es $g\cdot f$, pues $\varphi(x+n)=\varphi(x)$ para $n\in\mathbb{Z}$. Por lo tanto el punto final de ese levantamiento es $\Phi([g][f])$, y se tiene
\[
  \Phi([g][f])=(\tilde g\cdot f')(1)=n+m=\Phi[g]+\Phi[f].
\]
Esto dice que $\Phi$ es un homomorfismo.

\medskip\noindent\textit{$\Phi$ es suprayectiva.}
Sea $n\in\mathbb{Z}$. Se define $\tilde f\colon I\to\mathbb{R}$ por $\tilde f(s)=ns$, el camino que va de $0$ a $n$, y se proyecta al círculo: $f=\varphi\circ\tilde f\colon I\to S^1$. Este es un lazo basado en $1$ (pues $f(0)=1$ y $f(1)=e^{2\pi i n}=1$). Como el levantamiento es único, $\tilde f$ es el levantamiento de $f$ que comienza en $0$, así que $\Phi[f]=\tilde f(1)=n$.

\medskip\noindent\textit{$\Phi$ es inyectiva.}
Para esto basta ver que el núcleo es trivial. Sea $[f]\in\pi_1(S^1,1)$ con $\Phi[f]=0$. Entonces $\tilde f(1)=0$, y como $\tilde f(0)=0$, el camino $\tilde f$ es un lazo en $\mathbb{R}$ basado en $0$. Como ejercicio sencillo, $\tilde f$ es homotópico rel $\{0,1\}$ al camino constante en $0$ (por ejemplo, con la homotopía lineal). Proyectando esa homotopía a $S^1$ con $\varphi$ se obtiene una homotopía relativa a los extremos de $f$ al lazo constante en $1$. Es decir, $f$ se puede deformar rel $\{0,1\}$ al lazo constante, y por lo tanto $[f]$ es el neutro de $\pi_1(S^1,1)$. Así el núcleo es trivial y $\Phi$ es inyectiva.

Como $\Phi$ es un homomorfismo biyectivo, es un isomorfismo, y esto demuestra el teorema.
\end{proof}

Este es el primer cálculo explícito de un grupo fundamental y, de manera sorprendente, no fue tan complicado. Tiene varias aplicaciones, y en los siguientes videos se verán tres de ellas.

\end{document}
